\documentclass[../main.tex]{subfiles}

\begin{document}

\section{GENIE} \label{sec:genie}
  GENIE~\cite{Andreopoulos_2010}
  (\textit{Generate Events for Neutrino Interaction Experiments}) è un
  generatore di eventi Monte Carlo ampiamente utilizzato dalla comunità
  sperimentale della fisica dei neutrini.
  È un sistema software su larga scala basato su ROOT sviluppato
  interamente in \CPP.
  GENIE fornisce interazioni in un ampio intervallo energetico, da $\sim$1~MeV
  a $\sim$1~PeV, per tutti i bersagli nucleari e tutti i sapori di neutrini,
  include numerosi modelli fisici che possono essere classificati in: modelli di
  fisica nucleare, di sezione d'urto, di adronizzazione e di interazione nello
  stato finale (FSI). Inoltre, contiene diversi canali oltre il Modello Standard
  (BSM) e modelli ottimizzati (\textit{tuned}) derivanti dall'analisi dei dati
  globali.

  GENIE permette di personalizzare l'interazione dei neutrini con descrizioni
  dettagliate del flusso e della geometria del rivelatore.
  Le descrizioni del flusso sono fornite da simulazioni della
  \textit{beam-line} specifiche per l'esperimento, mentre quelle della
  geometria del rivelatore derivano dai disegni ingegneristici.\\
  Nel contesto della collaborazione DUNE, le simulazioni della
  \textit{beam-line} sono generate dal gruppo di simulazione del fascio
  LBNF (\textit{Long Baseline Neutrino Facility}).

\end{document}